\section{The marked map into the cusp neighborhood}

Retain the toric family, units, logarithms, section, and lattice
of Sections~1--6. Restrict to the compact tube $N_c$ over $|q|\leq\epsilon$.
Its boundary is smooth because the punctured family is a submersion.
Write $a=L(\theta,\beta)^t$ in the common $e$ marking.
For $q(v)=\epsilon e^{2\pi iv}$, use the logarithm
$\log q(v)=\log\epsilon+2\pi iv$ and set
\begin{equation}\label{ga:eq:j}
j_\infty[a,v]=
\left[\exp\bigl(2\pi i(\theta+\Omega(q(v))\beta)\bigr),q(v)\right]
\in\partial N_c.
\end{equation}
Exponentiation is coordinatewise. Brackets on the right mean the
quotient by the two specified periods.
Changing $\theta$ by an integer has no effect.
Changing $\beta$ by an integer multiplies by the corresponding periods.
At $v=1$, the matrix $\Omega$ is its value at zero plus $I_2$.
Thus \eqref{ga:eq:j} respects $[a,1]=[T_\infty a,0]$.
The same identity holds for derivatives across this seam.
The real period matrix is invertible, so this is a smooth diffeomorphism.
It carries the marked meridian to the specified section over the circle.

Give $F$ its complex orientation, which agrees with the $e$ orientation.
Give the disk its complex orientation and each boundary its outward-first
orientation. The boundary orientation of $N_c$ is fibre-first followed
by the positive meridian. The boundary of a punctured base has the
opposite meridian orientation. Thus \eqref{ga:eq:j} has the required
orientation reversal for attachment to the flat bundle over that base.

On integral normalized singular cubical chains define the actual map
\begin{equation}\label{ga:eq:kappa}
\kappa_\infty:C_*(B_{T_\infty})\longrightarrow C_*(N_c),
\qquad \sigma\longmapsto j_\infty\circ\sigma.
\end{equation}
This fixes the carrier and the map in every degree.
Finite cellular models of the compact pairs give the corresponding
finite chain map after cellular approximation and subdivision.

We relate this map to the toric cells without losing its meridian.
Use $L^{-1}$ to express the boundary in the $(\alpha,\beta)$ marking.
For $C=C_*^{\mathrm{cub}}(F;\mathbb Z)$ put
\[
R_n=C_n\oplus C_{n-1},\qquad
d_R(a,b)=(da+(T_*-I)b,-db).
\]
The interval-first prism comparison identifies $R$ with chains of $B_T$.
The toric forgetful map is represented by
\begin{equation}\label{ga:eq:K}
K:R_n\longrightarrow C_n^{\mathrm{cub}}(Z;\mathbb Z),
\qquad K(a,b)=p_*a+Hb,
\qquad (H\sigma)(v,t)=p_v(\sigma(t)).
\end{equation}
Here $p_v(\theta,\beta)=[y,\theta+vQ^{-1}y]$,
where $y$ represents $Q\beta$ in $H$.
The identity $dH+Hd=p_*T_*-p_*$ proves $dK=Kd_R$.

Let $j_r$ be the outer boundary identification from radial polar transport.
The self-map $j_r^{-1}j_\infty$ has the identity action on the full
based fundamental group, by the phase, deck, and section marking.
Choose equivariant lifts to the affine universal cover of $B_T$.
Their straight-line interpolation supplies a boundary homotopy.
Following that homotopy by the radial retraction gives a chain homotopy
\begin{equation}\label{ga:eq:comparison}
r_*\kappa_\infty-K=dS+Sd_R,
\end{equation}
with the prism comparison understood on the source.
The homotopy is fixed on the section.
It is defined by actual maps, rather than by an equality of homology ranks.
For nontrivial units, the smooth unit transport carries this construction
to the original family and its continuously specified marking.

The fibre part of \eqref{ga:eq:K} in degrees two and three is
\begin{equation}\label{ga:eq:p23}
P_2=\begin{pmatrix}
0&-1&1&1&-1&0\\0&1&0&0&0&0\\0&0&0&0&1&0\\0&0&0&0&0&1
\end{pmatrix},\qquad
P_3=\begin{pmatrix}0&0&1&0\\0&0&0&1\end{pmatrix}.
\end{equation}
Their source bases are the ordered exterior bases of
$(\alpha_1,\alpha_2,\beta_1,\beta_2)$.
Their target bases are $(E_0,E_1,E_2,h)$ and
$(h\otimes a_1,h\otimes a_2)$.
In the $e$ basis these matrices become $P_n\bigwedge^n L^{-1}$.
The term $H$ in \eqref{ga:eq:K} remains part of the boundary map.

\section{Integral meridian cycles}

Put $A_n=\bigwedge^n\mathbb Z^4$ and $N_n=\bigwedge^nT-I$.
Set $A_n=0$ outside $0\leq n\leq4$.
Write $\varepsilon_1,\varepsilon_2$ for the standard vectors in
the angular and polygon coordinate lattices $\mathbb Z^2$.
The mapping-torus chain sequence gives
\begin{equation}\label{ga:eq:wang}
0\longrightarrow A_n/N_nA_n\longrightarrow H_n(B_T;\mathbb Z)
\mathrel{\mathop{\longrightarrow}^{\pi}}\ker N_{n-1}\longrightarrow0.
\end{equation}
The left map is inclusion of the fibre.
Section~5 identifies its source with $H_n(Z;\mathbb Z)$ through $p_*$.
We now give an integral splitting for which $k_*$ vanishes on the right summand.

For a $T$-invariant oriented subtorus $V\subset F$, let $M(V)$ be its
mapping torus inside $B_T$. Orient it by the positive interval first,
followed by the stated basis of $V$.
Its fundamental class projects under $\pi$ to $[V]$.
All exterior products below are in the displayed order.
The following subtorus classes give integral bases of $\ker N_r$:
\begin{equation}\label{ga:eq:invariants}
\begin{array}{c|l}
r&[V]\\\hline
0&1\\
1&\alpha_1,\ \alpha_2\\
2&\alpha_1\alpha_2,\ \alpha_1\beta_1,\ \alpha_2\beta_2,
 (\alpha_1+\alpha_2)(\beta_1+\beta_2)\\
3&\alpha_1\alpha_2\beta_1,\ \alpha_1\alpha_2\beta_2\\
4&\alpha_1\alpha_2\beta_1\beta_2.
\end{array}
\end{equation}
Here juxtaposition denotes an exterior product.
For $r=0$, $M(V)$ is the marked meridian.
Each displayed sublattice is primitive and $T$-stable.
The restricted monodromy has determinant one.

To check the degree-two basis, expand $T\beta_j=\beta_j+\alpha_j$.
An invariant two-vector has no $\beta_1\beta_2$ term, and its coefficients
of $\alpha_1\beta_2$ and $\alpha_2\beta_1$ are equal.
The four displayed vectors form an integral basis of these solutions.
In degree three the only invariant vectors are integral combinations of
$\alpha_1\alpha_2\beta_1$ and $\alpha_1\alpha_2\beta_2$.
The remaining degrees follow directly from the same expansion.

Each class $[M(V)]$ in \eqref{ga:eq:invariants} maps to zero under $k_*$.
For purely angular $V$, the map is independent of the meridian coordinate.
For $V=\langle\alpha_j,\beta_j\rangle$, its base lies on the closed
line $\beta=t\varepsilon_j$. Representatives in $H$ have
$Q^{-1}y=(t-m)\varepsilon_j$ for an integer $m$ on each segment.
The angular shift $v(t-m)\varepsilon_j$ remains in the angular direction of $V$.
Thus its image lies in a two-dimensional quotient of rectangles times circles.
For $V=\langle\alpha_1+\alpha_2,\beta_1+\beta_2\rangle$, the same
argument uses the line $t(\varepsilon_1+\varepsilon_2)$ and seam translations
by $\varepsilon_1+\varepsilon_2$.
Its image also has a finite CW description of dimension at most two.
The three-dimensional subtori in the table have the entire angular torus
and a one-dimensional base. Their images have dimension at most three.
Finally, $k$ maps the five-dimensional class $[B_T]$ into the four-dimensional $Z$.
These factorizations prove the asserted vanishing in every degree.

\begin{theorem}\label{ga:thm:boundary}
The boundary homology is free, with ranks $(1,3,6,6,3,1)$.
For $0\leq n\leq4$, the map induced by the actual cusp attachment is
\begin{equation}\label{ga:eq:split}
H_n(B_T;\mathbb Z)=H_n(Z;\mathbb Z)\oplus\ker N_{n-1}
\mathrel{\mathop{\longrightarrow}^{k_*}}H_n(Z;\mathbb Z),
\qquad (z,u)\longmapsto z.
\end{equation}
The right summand is represented by the mapping-torus cycles above.
In degree five the map is zero.
\end{theorem}
\begin{proof}
The fibre quotient in \eqref{ga:eq:wang} maps isomorphically to $H_n(Z)$.
The mapping-torus cycles give a section of $\pi$, since their projected
classes form the integral basis \eqref{ga:eq:invariants}.
Their images under $k_*$ vanish by the explicit factorizations.
This proves \eqref{ga:eq:split}. Both summands are free.
Their ranks give the stated boundary ranks.
Equation~\eqref{ga:eq:comparison} identifies this result with
$r_*\kappa_{\infty*}$ in the original units and marking.
\end{proof}

In particular, the degree-two kernel has the actual torus basis
\begin{equation}\label{ga:eq:kernel2}
z_j=[M(\langle\alpha_j\rangle)],\qquad j=1,2.
\end{equation}
Explicitly, their fundamental squares are $(v,s)\mapsto[s\alpha_j,v]$.
Their boundaries cancel because $T\alpha_j=\alpha_j$.
The cusp retraction sends each square to $[0,s\varepsilon_j]$, independently of $v$.
The degree-three kernel has the four mapping-torus classes from row $r=2$.
These specify the meridian contributions as well as the fibre specialization.

\section{The integral local attachment cone}

For a chain map $f:B\to A$, use
$\operatorname{Cone}(f)_n=A_n\oplus B_{n-1}$ and
$d(a,b)=(da+fb,-db)$.
For the boundary inclusion \eqref{ga:eq:kappa}, this cone computes
$H_*(N_c,\partial N_c;\mathbb Z)$.
Indeed, inclusion of a boundary subcomplex is injective on cellular chains.
The cone maps onto the relative quotient, and its kernel is an identity cone.
Its contraction is $(c,b)\mapsto(0,c)$ in the identity-cone coordinates.

We describe the integral reduction used in the calculation.
A finite free chain complex with free homology splits as its homology
with zero differential plus contractible pairs.
To prove this, split the two exact sequences
\[
0\to Z_n\to C_n\to B_{n-1}\to0,
\qquad 0\to B_n\to Z_n\to H_n\to0.
\]
Here $Z_n=\ker d_n$ and $B_n=\operatorname{im}d_{n+1}$ in the complex $C$.
The quotients are free, so both sequences split over $\mathbb Z$.
The differential on the remaining pairs is the identity.
These splittings give maps $i,p,h$ with $pi=I$ and $I-ip=dh+hd$.
For a chain map $f:B\to A$, put $f_H=p_Afi_B$.
Then
\[
f-i_Af_Hp_B=dU+Ud,\qquad U=h_Af+i_Ap_Afh_B.
\]
Mapping cones of homotopic maps are isomorphic by their triangular
change of coordinates. Thus this operation preserves integral cone homology.

Both complexes here have free homology by Theorem~\ref{ga:thm:boundary}
and the local specialization theorem.
Consequently the actual local cone reduces to a finite free complex $D$.
Its degree-zero through degree-six ranks are
\begin{equation}\label{ga:eq:ranks}
(1,3,7,8,7,3,1).
\end{equation}
Write $r_n=(1,2,4,2,1)$ for the target ranks and
$b_n=(1,3,6,6,3,1)$ for the boundary ranks, with zero outside these ranges.
In the ordered coordinates $D_n=\mathbb Z^{r_n}\oplus\mathbb Z^{b_{n-1}}$,
the only nonzero block of $d_n$ is the upper-right block
\begin{equation}\label{ga:eq:reduced-d}
\begin{pmatrix}I_{r_{n-1}}&0\end{pmatrix}
:\mathbb Z^{b_{n-1}}\longrightarrow\mathbb Z^{r_{n-1}},
\qquad 1\leq n\leq5.
\end{equation}
The differential $d_6$ vanishes.
Thus the nonzero Smith factors of $d_1,\ldots,d_5$ are respectively
$1$, $(1,1)$, $(1,1,1,1)$, $(1,1)$, and $1$.

\begin{theorem}\label{ga:thm:relative}
The actual cusp boundary map has integral cone homology
\[
H_n(\operatorname{Cone}\kappa_\infty)=
\begin{cases}
0&n=0,1,\\
\mathbb Z&n=2,6,\\
\mathbb Z^2&n=3,5,\\
\mathbb Z^4&n=4,\\
0&\text{otherwise}.
\end{cases}
\]
\end{theorem}
\begin{proof}
In each degree the outgoing differential selects the first $r_{n-1}$
coordinates of the shifted boundary summand.
Its kernel is the full target summand and the remaining shifted coordinates.
The incoming differential identifies its domain coordinates bijectively
with the entire target summand. Therefore the quotient of cycles by
boundaries is freely generated by those remaining shifted coordinates.
Their ranks are $(0,0,1,2,4,2,1)$.
Every elimination uses a unit coefficient, so the quotient has no torsion.
The preceding chain contractions identify this exact integer calculation
with the cone of \eqref{ga:eq:kappa}.
\end{proof}

\section{The first homology of the common attachment}

We specify the remaining real topological data.
Let $\Sigma$ be an oriented sphere with three open disks removed.
Its marked positive meridians satisfy $\gamma_0\gamma_1\gamma_\infty=1$,
with right-to-left path multiplication. The flat torus bundle has monodromies
\[
T_0=\begin{pmatrix}0&-1&-1&0\\1&-1&0&0\\0&0&1&0\\0&0&0&1\end{pmatrix},\qquad
T_1=\begin{pmatrix}-1&1&-1&2\\-1&0&-1&1\\1&1&2&-1\\0&0&0&1\end{pmatrix}.
\]
They satisfy $T_0T_1T_\infty=I$ in the common $e$ basis.

Here are the finite fillings and their markings.
Set $u_0=e_3$, $u_1=-e_1+e_3$, $u_2=-e_1-e_2+e_3$.
Let $X_0$ be their three-torus and let $P_0$ cyclically permute this basis.
Put $U_0=u_0+u_1+u_2=\delta$, where $\delta=-2e_1-e_2+3e_3$.
For the other three-torus $X_1$ use the $f_1,f_2,f_3$ basis and put
\[
P_1=\begin{pmatrix}0&-1&-1\\1&0&0\\0&0&1\end{pmatrix},\quad
U_1=(-1,-1,2)^t,\quad
C=\begin{pmatrix}1&1&0&1\\0&1&0&0\\0&-1&1&0\\0&0&0&1\end{pmatrix}.
\]
The matrix $C$ sends the $f$ marking to the $e$ marking.
For $\rho=e^{2\pi i/3}$ define
\[
\begin{aligned}
N_0&=(X_0\times S^1_t\times D_s)/
 (x,t,s)\sim(P_0^{-1}x+U_0/3,t+1/3,\rho s),\\
N_1&=(X_1\times S^1_t\times D_s)/
 (x,t,s)\sim(P_1x-U_1/4,t-1/4,is).
\end{aligned}
\]
The finite actions are free because their nonidentity powers translate $t$.
The boundary maps are
\[
j_0[x,t,v]=[x-2vU_0/3,t+v/3,e^{2\pi iv/3}],\qquad
j_1[x,t,v]=[x-vU_1/4,t-v/4,e^{2\pi iv/4}].
\]
For $j_1$, first apply $C^{-1}$ to the fibre marking.
Substitution in these quotients verifies the mapping-torus endpoint relations.
Their core retractions, with $[y,\tau+1]=[P_i y,\tau]$, are
\[
[x,t,v]\longmapsto[x+2tU_0,3t+v],\qquad
[x,t,v]\longmapsto[x-tU_1,4t-v].
\]
Let $W$ be the flat bundle with these two fillings attached.
Its remaining boundary is $B_{T_\infty}$.
This defines a compact real manifold without assuming a global complex family.

Choose compatible finite chains $A=C_*(E)$, $L_i=C_*(N_i)$ and
$B_i=C_*(B_{T_i})$. Their common attachment complex is
\[
\mathcal W=\operatorname{Cone}\bigl(
(b_0,b_1)\mapsto(\iota_0b_0+\iota_1b_1,-j_{0*}b_0,-j_{1*}b_1)\bigr).
\]
The cusp boundary map into it is
$\beta(b)=(b_{\infty*}b,0,0;0,0)$.
Use the same boundary chain complex on both arrows of the attachment.
Subdivision and the cellular-to-singular comparisons transport the
cubical formulas and their specified homotopies to this finite diagram.
The actual attachment has the chain model
\begin{equation}\label{ga:eq:global}
\operatorname{Cone}\left((\beta,-\kappa_\infty):
C_*(B_{T_\infty})\to\mathcal W\oplus C_*(N_c)\right).
\end{equation}
The ordinary union comparison has contractible identity-cone kernels.
It therefore identifies \eqref{ga:eq:global} with the integral chains
of the specified smooth union $Y_{\mathrm{top}}=W\cup_{j_\infty}N_c$.
The retraction and \eqref{ga:eq:comparison} replace its second map by $-K$.
For the homotopy in that equation, the final change of coordinates is
$(w,z,b)\mapsto(w,z-Sb,b)$.
Substitution gives a chain isomorphism between the two core cones.

The finite core maps give the fundamental-group relations
\[
\pi_1(W)=\langle\Lambda,\gamma_0,\gamma_1\mid
[\Lambda,\Lambda]=1,\ \gamma_i\lambda\gamma_i^{-1}=T_i\lambda,
\ \gamma_0^3=w_0,\ \gamma_1^4=w_1\rangle,
\]
where $w_0=(4,2,-6,1)^t$ and $w_1=(1,1,-3,-1)^t$.
Indeed the core maps are coverings in the fourth period, of degrees three
and four, and the positive meridians map to the positive and negative core loops.
Their power relations are $w_0=e_4-2\delta$ and $w_1=-Cf_4-\delta$.
Van Kampen gives the displayed presentation.
The images of $T_0-I$ and $T_1-I$ generate
$\langle e_1,e_2,e_3\rangle$ integrally.
Abelianization therefore reduces to $3x_0=c$ and $4x_1=-c$.
The relation matrix
\[
\begin{pmatrix}-1&3&0\\1&0&4\end{pmatrix}
\]
has Smith factors $(1,1)$ and primitive kernel vector $(12,4,-3)^t$.
Thus $H_1(W)=\mathbb Zg$, with $e_4=12g$ and $\gamma_\infty=-g$.

The ordered basis of $H_1(B_{T_\infty})$ is
$(\beta_1,\beta_2,\gamma_\infty)$, where $\beta_1=e_3$ and $\beta_2=e_4$.
For the target basis $(g,\beta_1,\beta_2)$, the actual degree-one
attachment matrix is
\begin{equation}\label{ga:eq:global1}
(b_{\infty*},-k_*)_1=
\begin{pmatrix}0&12&-1\\-1&0&0\\0&-1&0\end{pmatrix}.
\end{equation}
Its determinant is $-1$. It is an isomorphism over $\mathbb Z$.
The degree-zero map is $(1,-1)^t$.
The cone exact sequence now gives
\begin{equation}\label{ga:eq:first}
H_0(Y_{\mathrm{top}};\mathbb Z)=\mathbb Z,
\qquad H_1(Y_{\mathrm{top}};\mathbb Z)=0.
\end{equation}

There is also a precise description of the next required map.
The cusp map is surjective on $H_2$, and its kernel is
$\mathbb Zz_1\oplus\mathbb Zz_2$ from \eqref{ga:eq:kernel2}.
Each class in $H_2(Z)$ has an integral fibre lift, and hence an integral
lift to $H_2(B_T)$. Eliminating that target summand in the cone exact
sequence gives
\begin{equation}\label{ga:eq:next}
H_2(Y_{\mathrm{top}};\mathbb Z)
\cong H_2(W;\mathbb Z)/
\langle b_{\infty*}z_1,b_{\infty*}z_2\rangle.
\end{equation}
The classes on the right are the actual cusp-meridian tori, rather than
two unnamed elements with prescribed ranks.
Their inclusion into the common complex must be evaluated to determine
this group, including any torsion.

For application to a global complex family, a further geometric comparison
must identify its punctured torus bundle with $E$ relative to all three
marked boundaries and the section. At each finite puncture, the lifted
translation on the reduced divisor torsor must extend smoothly, commute
with the cyclic deck action, and have a proper quotient with the displayed periods.
The retained section must realize the conormal comparisons
$e=s^{-2}e'$ at the order-three puncture and $e=s^{-1}e'$ at the order-four puncture.
These are existence and compatibility conditions on maps of the actual
torsors. They do not follow from their integral monodromy matrices.
The local cusp calculation and \eqref{ga:eq:first} leave those conditions
and the remaining homology of the common attachment distinct.
